% ============================================================================
% APPENDIX L: LIT-ACEMOGLU - Institutions, Automation & Political Economy
% Evidence-Based Framework for Economic and Social Behavior
% ============================================================================
% Category: LIT- (Literature Integration)
% Language: English
% Code: L
% Version: 2.0
% Last Updated: 2026-01-17
% Primary Dependencies: DOMAIN-HR, DOMAIN-GROWTH, CORE-WHEN, FORMAL-A
% EBF Integration: Institutions as fundamental Ψ that shapes all economic outcomes
% ============================================================================

\chapter*{Appendix L: Daron Acemoglu Research Integration}
\addcontentsline{toc}{chapter}{Appendix L: LIT-ACEMOGLU - Institutions, Automation \& Political Economy}
\label{app:acemoglu}

% ============================================================================
% CROSS-REFERENCE MAP
% ============================================================================
\begin{tcolorbox}[colback=blue!5!white,colframe=blue!75!black,title=Cross-Reference Map]
\textbf{Builds On:}
\begin{itemize}
    \item Appendix A (FORMAL-AXIOMS): Economic foundations
    \item Appendix V (CORE-WHEN): Context as fundamental driver
    \item Appendix O (LIT-AUTOR): Labor market task framework
\end{itemize}

\textbf{Supports:}
\begin{itemize}
    \item Appendix AJ (DOMAIN-HR): Labor market applications
    \item Appendix Z (DOMAIN-GROWTH): Growth theory applications
    \item Appendix RO (LIT-ROTH): Institutional design
\end{itemize}

\textbf{Related Literature:}
\begin{itemize}
    \item Appendix K (LIT-FEHR): Behavioral micro-foundations
    \item Appendix M (LIT-SHLEIFER): Behavioral finance institutions
    \item Appendix P (LIT-DUFLO): Development experiments
\end{itemize}
\end{tcolorbox}

% ============================================================================
% CHAPTER LINKAGE
% ============================================================================
\begin{tcolorbox}[colback=green!5!white,colframe=green!75!black,title=Chapter Linkage]
\textbf{Primary:} Chapter 9 (Context) -- Institutions as fundamental $\Psi$

\textbf{Secondary:}
\begin{itemize}
    \item Chapter 5 (Complementarity): $\gamma > 0$ between institutions and behavior
    \item Chapter 10.8 (Aggregate): Macro-level welfare effects
    \item Chapter 14 (Applications): Policy implications
\end{itemize}
\end{tcolorbox}

% ============================================================================
% ABSTRACT
% ============================================================================
\begin{abstract}
This appendix integrates Daron Acemoglu's research on institutions, automation, and political economy into the Evidence-Based Framework. Acemoglu's work demonstrates that institutions ($\Psi_{inst}$) are the fundamental cause of long-run prosperity differences, while his task-based automation framework shows how technology reshapes labor markets. We formalize six core axioms (AC-1 to AC-6) connecting institutional economics and automation theory to EBF. Key parameters include the automation displacement effect ($\delta_{robot} = -0.20\%$ per robot/1000 workers) and the automation share of wage inequality (50-70\%). Nobel Prize 2024 (with Simon Johnson and James Robinson).
\end{abstract}

% ============================================================================
% QUICK REFERENCE
% ============================================================================
\begin{tcolorbox}[colback=gray!5!white,colframe=gray!75!black,title=Quick Reference]
\textbf{Appendix:} L (LIT)

\textbf{Researcher:} Daron Acemoglu (MIT, Nobel 2024)

\textbf{Core Contributions:}
\begin{itemize}
    \item Institutions as fundamental cause of development
    \item Task-based framework for automation and labor
    \item Political economy of institutional change
    \item Directed technical change theory
\end{itemize}

\textbf{Key Parameters:}
\begin{itemize}
    \item Robot displacement: $\delta_{robot} = -0.20\%$ employment per robot/1000 workers
    \item Robot wage effect: $\delta_{wage} = -0.42\%$ per robot/1000 workers
    \item Automation → inequality: 50-70\% of US wage structure change
    \item Institutional persistence: Colonial origins explain 60\%+ of current income
\end{itemize}

\textbf{EBF Parameters Updated:}
\begin{itemize}
    \item $\Psi_{inst}$ (institutional context) $\rightarrow$ fundamental cause
    \item $\Psi_{tech}$ (technological context) $\rightarrow$ task displacement
    \item $\gamma_{inst-behavior}$ (institution-behavior complementarity)
\end{itemize}
\end{tcolorbox}

% ============================================================================
% FUNDAMENTAL QUESTION
% ============================================================================
% -----------------------------------------------------------------------------
% APPENDIX SCOPE BOX
% -----------------------------------------------------------------------------

\begin{tcolorbox}[
    colback=orange!5!white,
    colframe=orange!75!black,
    title={\textbf{Appendix Scope: Ziel / In-Scope / Out-of-Scope / Constraints}},
    fonttitle=\bfseries
]

\textbf{Ziel (Objective):}
\begin{itemize}[nosep]
    \item Why do some nations prosper while others remain poor, and how does technological change affect the distribution of prosperity?
\end{itemize}

\textbf{In-Scope:}
\begin{itemize}[nosep]
    \item Fundamental Question
    \item Core Axioms: Acemoglu's Contributions
    \item Key Empirical Results
    \item Integration with Evidence-Based Framework
\end{itemize}

\textbf{Out-of-Scope (delegiert an andere Appendices):}
\begin{itemize}[nosep]
    \item FORMAL-AXIOMS $\rightarrow$ Appendix A
    \item CORE-WHEN $\rightarrow$ Appendix V
    \item LIT-AUTOR $\rightarrow$ Appendix O
    \item DOMAIN-HR $\rightarrow$ Appendix AJ
    \item DOMAIN-GROWTH $\rightarrow$ Appendix Z
\end{itemize}

\textbf{Constraints (Anwendungsgrenzen):}
\begin{itemize}[nosep]
    \item [TODO: Define when this appendix does NOT apply]
    \item [TODO: Define required prerequisites]
    \item [TODO: Define known limitations]
\end{itemize}

\textbf{Lieferobjekte (Deliverables):}
\begin{itemize}[nosep]
    \item 2 Table(s)
    \item 9 Equation(s)
    \item 12 Axiom(s)
    \item Worked Example(s)
\end{itemize}

\end{tcolorbox}


\section{Fundamental Question}

\begin{tcolorbox}[colback=red!5!white,colframe=red!75!black,title=Central Question]
\textbf{Why do some nations prosper while others remain poor, and how does technological change affect the distribution of prosperity?}
\end{tcolorbox}

Acemoglu's research addresses two fundamental questions:

\begin{enumerate}
    \item \textbf{Institutions and Development}: Why do income differences across countries persist for centuries? Answer: Institutions determine incentives, and institutions persist.
    \item \textbf{Automation and Labor}: How does technology affect employment and wages? Answer: Technology displaces tasks, but new tasks can reinstate labor.
\end{enumerate}

For EBF, Acemoglu provides:
\begin{itemize}
    \item The theoretical foundation for $\Psi_{inst}$ as fundamental context
    \item Quantified parameters for automation effects on labor markets
    \item A framework connecting macro institutions to micro behavior
\end{itemize}

% ============================================================================
% THEORETICAL FOUNDATIONS
% ============================================================================
\section{Theoretical Foundations}

\subsection{Institutions as Fundamental Cause}

Acemoglu's institutional theory can be summarized:

\begin{equation}
\Psi_{inst} \xrightarrow{\text{incentives}} \text{Investment, Innovation} \xrightarrow{\text{accumulation}} \text{Prosperity}
\end{equation}

\textbf{Key distinction}:
\begin{itemize}
    \item \textbf{Inclusive institutions}: Secure property rights, rule of law, open markets $\rightarrow$ growth
    \item \textbf{Extractive institutions}: Elite capture, weak enforcement $\rightarrow$ stagnation
\end{itemize}

\textbf{Persistence mechanism}: Colonial settlers chose institutions based on mortality rates. High mortality $\rightarrow$ extractive institutions $\rightarrow$ persist today.

\subsection{Task-Based Framework for Automation}

The task framework (with Autor, Restrepo) reconceptualizes labor markets:

\begin{equation}
Y = F\left(\int_0^1 y(i) di\right) \text{ where } y(i) = A_L(i) \cdot l(i) + A_K(i) \cdot k(i)
\end{equation}

\textbf{Key insight}: Technology doesn't affect ``labor'' directly---it affects \textit{tasks}:
\begin{itemize}
    \item \textbf{Displacement effect}: Automation replaces labor in existing tasks
    \item \textbf{Reinstatement effect}: New tasks create demand for labor
    \item \textbf{Net effect}: Depends on balance of these forces
\end{itemize}

% ============================================================================
% CORE AXIOMS
% ============================================================================
\section{Core Axioms: Acemoglu's Contributions}

We formalize six axioms capturing Acemoglu's key contributions:

\subsection{Axiom AC-1: Institutional Primacy}

\begin{tcolorbox}[colback=blue!5!white,colframe=blue!75!black,title=Axiom AC-1: Institutional Primacy]
\textbf{Statement}: Institutions are the fundamental cause of long-run economic development. Geography, culture, and other factors operate primarily through their effect on institutions.

\textbf{Formal Definition}:
\begin{equation}
\frac{\partial \ln Y_{LR}}{\partial \Psi_{inst}} > \frac{\partial \ln Y_{LR}}{\partial \Psi_{geo}}, \frac{\partial \ln Y_{LR}}{\partial \Psi_{culture}}
\end{equation}

\textbf{Empirical Support}:
\begin{itemize}
    \item Colonial Origins (2001): Settler mortality instruments for institutions, R² > 0.6
    \item Reversal of Fortune: Pre-1500 rich areas now poor (extractive institutions)
    \item Korea comparison: Same geography/culture, different institutions $\rightarrow$ 10x income gap
\end{itemize}

\textbf{Evidence Tier}: 2 (Natural experiment with instrumental variables)
\end{tcolorbox}

\subsection{Axiom AC-2: Institutional Persistence}

\begin{tcolorbox}[colback=blue!5!white,colframe=blue!75!black,title=Axiom AC-2: Institutional Persistence]
\textbf{Statement}: Institutions persist over long periods through self-reinforcing mechanisms, even when de jure rules change.

\textbf{Formal Definition}:
\begin{equation}
\text{Corr}(\Psi_{inst,t}, \Psi_{inst,t+\tau}) > 0.5 \text{ for } \tau > 100 \text{ years}
\end{equation}

\textbf{Mechanisms}:
\begin{itemize}
    \item De facto power persists even when de jure power changes
    \item Elites adapt to institutional changes
    \item Complementarities between institutions create lock-in
\end{itemize}

\textbf{Empirical Support}:
\begin{itemize}
    \item Persistence of Power (2008): Elites maintain influence after democratization
    \item Colonial institutions: 400+ years of persistence
\end{itemize}

\textbf{Evidence Tier}: 2 (Historical natural experiments)
\end{tcolorbox}

\subsection{Axiom AC-3: Task Displacement}

\begin{tcolorbox}[colback=blue!5!white,colframe=blue!75!black,title=Axiom AC-3: Task Displacement Effect]
\textbf{Statement}: Automation reduces labor demand by displacing workers from tasks they previously performed. The effect is concentrated in routine tasks.

\textbf{Formal Definition}:
\begin{equation}
\frac{\partial L}{\partial A_{robot}} = -\delta_{robot} < 0
\end{equation}

where $\delta_{robot} = 0.20\%$ employment reduction per robot per 1000 workers.

\textbf{Empirical Support}:
\begin{itemize}
    \item Robots and Jobs (2020): $\delta_{robot} = 0.20\%$, $\delta_{wage} = 0.42\%$
    \item Concentrated in manufacturing, routine tasks
    \item Commuting zone analysis with Bartik-style instruments
\end{itemize}

\textbf{Evidence Tier}: 1 (Causal identification with instruments)
\end{tcolorbox}

\subsection{Axiom AC-4: Task Reinstatement}

\begin{tcolorbox}[colback=blue!5!white,colframe=blue!75!black,title=Axiom AC-4: Task Reinstatement Effect]
\textbf{Statement}: Technological change creates new tasks where labor has comparative advantage, potentially offsetting displacement effects.

\textbf{Formal Definition}:
\begin{equation}
\frac{\partial L}{\partial \text{NewTasks}} = +\delta_{new} > 0
\end{equation}

\textbf{Historical Evidence}:
\begin{itemize}
    \item 1900-2000: Massive automation, yet labor share relatively stable
    \item New occupations absorb displaced workers over time
    \item Rate of new task creation has slowed since 1980
\end{itemize}

\textbf{Evidence Tier}: 3 (Historical analysis)
\end{tcolorbox}

\subsection{Axiom AC-5: Automation and Inequality}

\begin{tcolorbox}[colback=red!5!white,colframe=red!75!black,title=Axiom AC-5: Automation Drives Inequality]
\textbf{Statement}: Automation is the primary driver of rising wage inequality since 1980, accounting for 50-70\% of changes in the US wage structure.

\textbf{Formal Definition}:
\begin{equation}
\Delta \text{Gini}_{1980-2016} = 0.5 \cdot \Delta \text{Automation} + 0.3 \cdot \Delta \text{Other} + \epsilon
\end{equation}

\textbf{Empirical Support}:
\begin{itemize}
    \item Econometrica (2022): 50-70\% of wage structure changes
    \item Task displacement concentrated among non-college workers
    \item Explains polarization better than skill-biased technical change
\end{itemize}

\textbf{Evidence Tier}: 1 (Causal decomposition analysis)
\end{tcolorbox}

\subsection{Axiom AC-6: Directed Technical Change}

\begin{tcolorbox}[colback=green!5!white,colframe=green!75!black,title=Axiom AC-6: Endogenous Direction of Technology]
\textbf{Statement}: The direction of technical change responds to economic incentives. Technology develops to use abundant factors more intensively.

\textbf{Formal Definition}:
\begin{equation}
\frac{\partial A_L/A_K}{\partial (L/K)} > 0 \text{ (weak induced bias)}
\end{equation}

\textbf{Implications}:
\begin{itemize}
    \item Large labor supply $\rightarrow$ labor-augmenting innovation
    \item High wages $\rightarrow$ labor-saving automation
    \item Policy can redirect technological development
\end{itemize}

\textbf{Evidence Tier}: 5 (Theoretical with empirical support)
\end{tcolorbox}

% ============================================================================
% KEY EMPIRICAL RESULTS
% ============================================================================
\section{Key Empirical Results}

\subsection{Parameter Estimates from Acemoglu's Research}

\begin{table}[h]
\centering
\caption{Key Parameters from Acemoglu's Research}
\label{tab:acemoglu-parameters}
\begin{tabular}{llcc}
\toprule
\textbf{Parameter} & \textbf{Context} & \textbf{Estimate} & \textbf{Source} \\
\midrule
\multicolumn{4}{l}{\textit{Automation Effects}} \\
Employment displacement & Robots & $-0.20\%$ per robot/1000 & JPE 2020 \\
Wage effect & Robots & $-0.42\%$ per robot/1000 & JPE 2020 \\
Inequality share & Automation & 50-70\% & Econometrica 2022 \\
\midrule
\multicolumn{4}{l}{\textit{Institutional Effects}} \\
Colonial origin R² & Income today & $> 0.60$ & AER 2001 \\
Institutional persistence & 400 years & $\rho > 0.50$ & Various \\
Inclusive vs extractive & Long-run growth & 2-3x difference & WNF 2012 \\
\bottomrule
\end{tabular}
\end{table}

\subsection{Worked Example: Robots and US Labor Markets}

\begin{tcolorbox}[colback=green!5!white,colframe=green!75!black,title=Worked Example: Robot Adoption 1990-2007]
\textbf{Setting}: US commuting zones, 1990-2007

\textbf{Data}:
\begin{itemize}
    \item Robot adoption: IFR data on industrial robots
    \item Employment: Census/ACS data
    \item Instrument: European robot adoption (industry-level)
\end{itemize}

\textbf{Results}:
\begin{itemize}
    \item One additional robot per 1000 workers:
    \begin{itemize}
        \item Employment-to-population: $-0.20$ percentage points
        \item Log wages: $-0.42\%$
    \end{itemize}
    \item Effects concentrated in:
    \begin{itemize}
        \item Manufacturing
        \item Routine manual tasks
        \item Non-college workers
    \end{itemize}
\end{itemize}

\textbf{EBF Interpretation}:
\begin{itemize}
    \item $\Psi_{tech}$ (robot adoption) shifts task allocation
    \item Displacement effect dominates in affected sectors
    \item $\gamma_{skill-task} < 0$: Routine skills become less complementary with technology
\end{itemize}
\end{tcolorbox}

% ============================================================================
% EBF INTEGRATION
% ============================================================================
\section{Integration with Evidence-Based Framework}

\subsection{Connection to 10C Framework}

\begin{table}[h]
\centering
\caption{Acemoglu's Contributions to 10C Framework}
\label{tab:acemoglu-9c}
\begin{tabular}{lll}
\toprule
\textbf{CORE} & \textbf{Acemoglu Contribution} & \textbf{Mechanism} \\
\midrule
CORE-WHO (AAA) & Levels of analysis & Individual, firm, nation \\
CORE-WHAT (C) & Prosperity dimensions & Income, employment, wages \\
CORE-HOW (B) & Institution-behavior $\gamma$ & Incentive complementarities \\
CORE-WHEN (V) & $\Psi_{inst}$, $\Psi_{tech}$ & Fundamental context \\
CORE-WHERE (BBB) & Quantified parameters & $\delta_{robot}$, persistence $\rho$ \\
CORE-STAGE (AW) & Development trajectory & Path dependence \\
\bottomrule
\end{tabular}
\end{table}

\subsection{The Institutions-Behavior Link}

Acemoglu's work provides the macro-to-micro link in EBF:

\begin{equation}
\Psi_{inst} \xrightarrow{\text{incentives}} \gamma_{ij} \xrightarrow{\text{behavior}} U^L_d(\Psi)
\end{equation}

\textbf{Example}: Secure property rights ($\Psi_{inst}$) increase returns to investment ($\gamma_{effort,return} > 0$), leading to more productive behavior.

\subsection{Policy Implications}

\begin{tcolorbox}[colback=yellow!5!white,colframe=yellow!75!black,title=Policy Implications from Acemoglu]
\begin{enumerate}
    \item \textbf{Institutions First}: Development requires institutional reform, not just aid
    \item \textbf{Technology Policy}: Direction of innovation can be influenced by policy
    \item \textbf{Automation Response}: New task creation requires active policy support
    \item \textbf{Inequality}: Automation effects require compensatory policies
\end{enumerate}

\textbf{EBF Design Principle}: Target $\Psi_{inst}$ for sustainable change, not just individual behavior.
\end{tcolorbox}

% ============================================================================
% CRITICAL FOUNDATIONS
% ============================================================================
\section{Critical Foundations}

\begin{tcolorbox}[colback=orange!5!white,colframe=orange!75!black,title=Critical Foundations]
\begin{enumerate}
    \item \textbf{North (1990)}: Institutions as rules of the game
    \item \textbf{Coase (1960)}: Property rights and transaction costs
    \item \textbf{Autor, Levy, Murnane (2003)}: Task framework origins
    \item \textbf{Romer (1990)}: Endogenous growth theory
    \item \textbf{Nunn (2009)}: Historical persistence mechanisms
\end{enumerate}
\end{tcolorbox}

% ============================================================================
% OPEN ISSUES
% ============================================================================
\section{Open Issues and Research Frontier}

\begin{enumerate}
    \item \textbf{AI vs. Previous Automation}: Will AI create enough new tasks?
    \item \textbf{Institutional Reform}: How to change extractive institutions?
    \item \textbf{Global Automation}: Developing country implications?
    \item \textbf{Political Economy of AI}: Who controls AI development direction?
    \item \textbf{Measurement}: Better metrics for institutional quality?
\end{enumerate}

% ============================================================================
% SUMMARY
% ============================================================================
\section{Summary}

\begin{tcolorbox}[colback=gray!5!white,colframe=gray!75!black,title=Chapter Summary]
Daron Acemoglu's research provides EBF with:

\textbf{Theoretical Contributions}:
\begin{itemize}
    \item Institutions as fundamental $\Psi$ for development
    \item Task-based framework for automation analysis
    \item Directed technical change theory
    \item Political economy of institutional persistence
\end{itemize}

\textbf{Quantified Parameters}:
\begin{itemize}
    \item Robot displacement: $-0.20\%$ employment per robot/1000
    \item Robot wage effect: $-0.42\%$ per robot/1000
    \item Automation share of inequality: 50-70\%
\end{itemize}

\textbf{EBF Integration}:
\begin{itemize}
    \item $\Psi_{inst}$ as fundamental context dimension
    \item $\Psi_{tech}$ as task-level technological context
    \item Institution-behavior complementarity $\gamma_{inst-behavior}$
\end{itemize}

\textbf{Key Insight}: For lasting change, reform institutions ($\Psi_{inst}$) rather than just individual behavior. Technology's effects depend on institutional context.
\end{tcolorbox}

% ============================================================================
% GLOSSARY
% ============================================================================
\section{Glossary}

Key terms defined in this appendix (see also Appendix G for complete glossary):

\begin{description}
    \item[Displacement Effect] Reduction in labor demand from automation of existing tasks
    \item[Extractive Institutions] Rules that concentrate power and wealth in elite hands
    \item[Inclusive Institutions] Rules that spread opportunity and enforce contracts broadly
    \item[Institutional Persistence] Tendency of institutions to remain stable over long periods
    \item[Reinstatement Effect] Creation of new tasks where labor has comparative advantage
    \item[Task] A unit of production that can be performed by labor or capital
    \item[Task-Based Framework] Approach analyzing labor markets through task allocation
\end{description}

% ============================================================================
% REFERENCES
% ============================================================================
\section{References}

\nocite{bcm_master}

\textbf{Primary Acemoglu Sources} (20 papers integrated):

\begin{itemize}
    \item Acemoglu, D., Johnson, S., \& Robinson, J.A. (2001). Colonial Origins. \textit{AER}.
    \item Acemoglu, D. (2002). Directed Technical Change. \textit{REStud}.
    \item Acemoglu, D. (2003). Labor- and Capital-Augmenting Technical Change. \textit{JEEA}.
    \item Acemoglu, D., Johnson, S., \& Robinson, J.A. (2005). Rise of Europe. \textit{AER}.
    \item Acemoglu, D., \& Robinson, J.A. (2006). Economic Origins of Dictatorship and Democracy. Cambridge.
    \item Acemoglu, D., Aghion, P., \& Zilibotti, F. (2006). Distance to Frontier. \textit{JEEA}.
    \item Acemoglu, D., \& Johnson, S. (2007). Disease and Development. \textit{JPE}.
    \item Acemoglu, D., \& Robinson, J.A. (2008). Persistence of Power. \textit{AER}.
    \item Acemoglu, D., \& Autor, D. (2011). Skills, Tasks, Technologies. \textit{Handbook of Labor Economics}.
    \item Acemoglu, D., \& Robinson, J.A. (2012). Why Nations Fail. Crown.
    \item Acemoglu, D., Garcia-Jimeno, C., \& Robinson, J.A. (2015). State Capacity. \textit{AER}.
    \item Acemoglu, D., Ozdaglar, A., \& Tahbaz-Salehi, A. (2016). Networks and Systemic Risk. \textit{Handbook}.
    \item Acemoglu, D., \& Restrepo, P. (2018). Race between Man and Machine. \textit{AER}.
    \item Acemoglu, D., \& Restrepo, P. (2019). Automation and New Tasks. \textit{JEP}.
    \item Acemoglu, D., \& Restrepo, P. (2020). Robots and Jobs. \textit{JPE}.
    \item Acemoglu, D. (2021). Harms of AI. \textit{NBER WP}.
    \item Acemoglu, D., \& Restrepo, P. (2022). Tasks, Automation, and Inequality. \textit{Econometrica}.
    \item Acemoglu, D., \& Johnson, S. (2023). Power and Progress. PublicAffairs. \citep{acemoglu2023power}
\end{itemize}

\end{document}
